\chapter{Galois Extensions}
\section{Fundamental Theorem of Galois Theory}
We begin this section by defining what it means for a field extension to be {\bf Galois}:
\begin{definition}
    A field extension $L:K$ is {\bf Galois} if $K=L^{\mathrm{Aut}_K(L)}$, the Galois group of a Galois extension is $\operatorname{Gal}(L/K):=\mathrm{Aut}_K(L)$.

    If $[L:K] < \infty$ (which we will assume to be the case for the rest of this course), then we may apply Theorem \ref{thm-galcondition} to get three equivalent characterizations of Galois extensions.
\end{definition}
\begin{example}
    \begin{enumerate}
        \item $\mathbb{Q}(\sqrt[3]{2}):\mathbb{Q}$ is not Galois, since it is not a normal extension by Example \ref{eg-normal}. Note that $\mathrm{Aut}_{\mathbb{Q}}(\mathbb{Q}(\sqrt[3]{2})) = \{id\}$ (Exercise) has only one element, which is strictly less than $[\mathbb{Q}(\sqrt[3]{2}):\mathbb{Q}] = 3$.
        \item  $\mathbb{Q}(i):\mathbb{Q}$ is a normal extension, since it is the splitting extension of $p(x) = x^2+1 \in \mathbb{Q}[x]$. Moreover, by Corollary \ref{cor-groupsize}, $|\operatorname{Gal}(\mathbb{Q}(i)/\mathbb{Q})| = |\mathrm{Aut}_{\mathbb{Q}}(\mathbb{Q}(i))| = [\mathbb{Q}(i):\mathbb{Q}] = 2$ has two elements, namely:
        $\operatorname{Gal}(\mathbb{Q}(i)/\mathbb{Q}) = \{id, \sigma\}$, where $id(a+bi):=a+bi$ and $\sigma(a+bi):=a-bi$ for $a+bi \in \mathbb{Q}(i)$.
    \end{enumerate}
\end{example}

We begin with an overview of Fundamental Theorem of Galois Theory: 
Suppose $L : K$ is Galois extension, the Fundamental Theorem of Galois Theory gives an 1-1 correspondence between 
\begin{center}
    $\{M\ |\ L \supseteq M \supseteq K\}$ and $\{H\ |\ H\leq \operatorname{Gal}(L /K)\}$. 
\end{center}
The relations between field extensions and Galois groups are shown in the following diagram: 

\[
\begin{tikzcd}
        & L &&&&& {\operatorname{Gal}(L/L)} \\
        {H_{1}} & {H_{2}} & {H_{3}} & {} & {} & {\operatorname{Gal}(L/H_{1})} & {\operatorname{Gal}(L/H_{2})} & {\operatorname{Gal}(L/H_{3})} \\
        & K &&&&& {\operatorname{Gal}(L/K)}
        \arrow[from=1-7, to=2-6]
        \arrow[from=1-7, to=2-7]
        \arrow[from=1-7, to=2-8]
        \arrow[from=2-1, to=1-2]
        \arrow[from=2-2, to=1-2]
        \arrow[from=2-3, to=1-2]
        \arrow[squiggly, tail reversed, from=2-5, to=2-4]
        \arrow[from=2-6, to=3-7]
        \arrow[from=2-7, to=3-7]
        \arrow[from=2-8, to=3-7]
        \arrow[from=3-2, to=2-1]
        \arrow[from=3-2, to=2-2]
        \arrow[from=3-2, to=2-3]
    \end{tikzcd}
\]

\begin{theorem}[Fundamental Theorem of Galois Theory, Part I]\label{thm-FTGT-1}
    Let $L:K$ be Galois extension, then there's an order reversing bijection between the sets
    \[
        \begin{tikzcd}
            {\{M\ |\ L\supseteq M \supseteq K\}} & {\{H\ |\ \{e\}\leq H\leq \operatorname{Gal}(L /K)\}}
            \arrow["\phi", shift left, from=1-1, to=1-2]
            \arrow["\varphi", shift left, from=1-2, to=1-1]
        \end{tikzcd}\]
given by
\begin{align*}
    \phi(M) &:= \operatorname{Gal}(L /M)\\
     L^{H} &=: \varphi(H)
\end{align*}
\end{theorem}

\begin{remark}
    \begin{enumerate}
        \item If $L:K$ is Galois, then for any $L\supseteq M\supseteq K$, $L : M$ is also Galois. To see why, $L:K$ is Galois implies $L:K$ is a splitting extension of some separable polynomial $p(x)\in K[x]$. Then the same polynomial is separable in $M[x]$ by Corollary \ref{cor-sepPolyFieldExt}. So $L$ is the splitting extension of separable extension $p(x)\in M[x]$. 
        \item For any $H\leq \operatorname{Gal}(L /K)$, if $\phi\in H$, then $\phi(k)=k$ for any $k\in K$. So $K\subseteq L^{H}$. The fact that $L^{H}$ is a field is left as an exercise. 
        \item By order reversing, we mean 
        \[
        (M_{1}\supseteq M_{2})\mapsto (\operatorname{Gal}(L /M_{1})\leq \operatorname{Gal}(L /M_{2}));
        \]
        \[
        (L^{H_{1}}\supseteq L^{H_{2}})\mapsfrom (H_{1}\subseteq H_{2}).
        \]
        \item Since $L:M$ is Galois, we have $L^{\mathrm{Aut}_{M}(L)}=M$ for all $L\supseteq M\supseteq K$. Therefore, $\varphi\circ \phi(M)=M$, i.e. $\varphi\circ \phi=id$. So we just need to show $\phi\circ\varphi=id$ in Theorem \ref{thm-FTGT-1}. 
    \end{enumerate}
\end{remark}

\begin{proof}[Proof of Theorem \ref{thm-FTGT-1}]
    We only show $H=\operatorname{Gal}(L /L^{H})$ for any $H\leq \operatorname{Gal}(L /K)$. We already have $H\leq \operatorname{Gal}(L /L^{H})$ since for any $\sigma \in H$, $\sigma:L\to L$ satisfies $\sigma(\gamma)=\gamma$ for any $\gamma\in L^{H}$. So $\sigma$ is an element in $\mathrm{Aut}_{L^{H}}(L)=\operatorname{Gal}(L /L^{H})$. 

    To see it is an equality, we make the following: 
    
    \begin{claim}
    Let $L : K$ be of finite degree (which is not necessarily Galois), and $H\leq \mathrm{Aut}_{K}(L)$. Then $[L:L^{H}]\leq \left| H \right| $.
    \end{claim}
    \begin{proof}[Proof of Claim]
        Suppose on the contrary, i.e. $[L:L^{H}]>\left| H \right| $. Let $n=[L:L^{H}]$ and $d=\left| H \right| $. Let $\{\sigma_{1},\dots,\sigma_{d}\}=H\leq \mathrm{Aut}_{K}(L)$,  and $\{\alpha_{1},\dots, \alpha_{n}\}$ be a basis of $L /L^{H}$. Set 
        \[
        v_{1}=\begin{pmatrix}
            \sigma_{1}(\alpha_{1})\\
            \sigma_{2}(\alpha_{1})\\
            \vdots\\
            \sigma_{d}(\alpha_{1})
        \end{pmatrix},
        v_{2}=\begin{pmatrix}
            \sigma_{1}(\alpha_{2})\\
            \sigma_{2}(\alpha_{2})\\
            \vdots\\
            \sigma_{d}(\alpha_{2})
        \end{pmatrix},
        \dots, 
        v_{n}=\begin{pmatrix}
            \sigma_{1}(\alpha_{n})\\
            \sigma_{2}(\alpha_{n})\\
            \vdots\\
            \sigma_{d}(\alpha_{n})
        \end{pmatrix}
        \in L^{d}
        \] 
        as $n$ vectors in the $d$-dimensional vector space $L^{d}$ over $L$. Since $n>d$, those vectors are linearly dependent over $L$. Assume $\sigma_{1}$ is the identity element in $H$. 

        Suppose ``$r\leq n$'' is the smallest integer such that there exists nontrivial linear relations between vectors above, and (after reordering the $v_i$'s if necessary) there exists $\beta_1, \dots, \beta_r \in L$ such that
        \[
        \beta_{1}v_{1}+\cdots \beta_{r}v_{r}=0. 
        \]
        WLOG, we may assume $\beta_{r}\neq 0$ here, then replace $\beta_{i}$ by $\beta_{i}'=\beta_{i}/ \beta_{r}$, one get 
        \[
        \beta_{1}'v_{1}+\cdots \beta_{r-1}'v_{r-1}+v_{r}=0. \tag*{(*)}
        \]
        for $\beta_1', \dots, \beta_{r-1}' \in L$. Pick the first row in (*) , we have 
        \[
        \beta_{1}'\sigma_{1}(\alpha_{1})+\cdots +\beta_{r-1}'\sigma_{1}(\alpha_{r-1})+\sigma_{1}(\alpha_{r})=0,
        \]
        which is 
        \[
        \beta_{1}'\alpha_{1}+\cdots \beta_{r-1}'\alpha_{r-1}+\alpha_{r}=0. 
        \]
        
        Since $\{\alpha_{1},\dots, \alpha_{n}\}$ forms a basis of $L$ as an $L^{H}$ vector space, they are $L^{H}$-linearly independent. So at least one $\beta_{i}'$ is not in $L^{H}$. WLOG, assume $\beta_{1}'\notin L^{H}$, that means there exists an element $\tau\in H$ such that $\tau(\beta_{1}') \neq \beta_{1}'$. 

        Pick the $i$-th row in (*)
        \[
        \beta_{1}'\sigma_{i}(\alpha_{1})+\cdots +\beta_{r-1}'\sigma_{i}(\alpha_{r-1})+\sigma_{i}(\alpha_{r})=0,
        \]
        apply $\tau$ to the equation, 
        \[
        \tau(\beta_{1}')\tau\sigma_{i}(\alpha_{1})+\cdots+\tau(\beta_{r-1}')\tau\sigma_{i}(\alpha_{r-1})+\tau\sigma_{i}(\alpha_{r})=0. 
        \]
        Note $\{\tau\sigma_{1}, \dots, \tau\sigma_{d}\} = \{\sigma_{1}, \dots, \sigma_{d}\} = H$, so one has
        \[
        \tau(\beta_{1}')\sigma_{j}(\alpha_{1})+\cdots+\tau(\beta_{r-1}')\sigma_{j}(\alpha_{r-1})+\sigma_{j}(\alpha_{r})=0.
        \]
        for all $1 \leq j \leq d$. Combine all the rows together, 
        \[
        \tau(\beta_{1}')v_{1}+\cdots+\tau(\beta_{r-1}')v_{r-1}+v_{r}=0.
        \]
    Subtract the above equation by $(*)$, one has:
        \[(\tau(\beta_{1}')-\beta_{1}')v_{1}+\cdots+(\tau(\beta_{r-1}')-\beta_{r-1}')v_{r-1}=0. \]
        with the first term $\tau(\beta_{1}')-\beta_{1}' \neq 0$. In other words, we find a smaller nontrivial linear relationship involving only ``$r-1\leq n$'' vectors, which contradict to the minimality of $r$. 
    \end{proof}

    The claim gives $[L:L^{H}]\leq |H|$. Since $H\leq \mathrm{Aut}_{L^H}(L)$, while Theorem \ref{thm-genKhomo} gives $|\mathrm{Aut}_{L^H}(L)|\leq [L:L^H]$, we have
    \[
    |H|\leq |\mathrm{Aut}_{L^H}(L)|\leq [L:L^H]\leq |H|.
    \]
    Hence $H=\mathrm{Aut}_{L^H}(L)=\operatorname{Gal}(L/L^H)$. 
\end{proof}

\begin{theorem}[Fundamental Theorem of Galois Theory, Part II]\label{thm-FTGT-2}
    Let $L:K$ be Galois extension, $M:K$ be a subextension of $L$. Then $M:K$ is Galois if and only if $\operatorname{Gal}(L /M)\lhd \operatorname{Gal}(L /K)$. In such case, $\operatorname{Gal}(M /K)\cong \operatorname{Gal}(L /K) /\operatorname{Gal}(L /M)$. 
\end{theorem}

\begin{proof}
    $(\Rightarrow)$: Suppose $M :K$ is Galois extension. Let $\gamma\in \operatorname{Gal}(L /K)$. We consider $\gamma|_{M}$: 
    \[
        \begin{tikzcd}
            M && L \\
            & K
            \arrow["{\gamma|_{M}}", from=1-1, to=1-3]
            \arrow[from=2-2, to=1-1]
            \arrow[from=2-2, to=1-3]
        \end{tikzcd}
    \]
    since $M :K$ is splitting field of some separable polynomial, by Theorem \ref{thm-splittingfieldsameimage} the image $\gamma|_{M}:M\to L$ are the same for all $\gamma\in \operatorname{Gal}(L /K)$. Pick the identity map, we have the result $\gamma|_{M}(M)=id_{M}(M)=M$, so $\gamma|_{M}\in \operatorname{Gal}(M/K)$. 

    Define the map
    \[
    \Phi: \operatorname{Gal}(L /K)\to \operatorname{Gal}(M /K)
    \]
    by \[ 
    \gamma\mapsto \gamma|_{M}.
    \] 
    $\Phi$ is clearly a group homomorphism. $\Phi(\gamma)=id$ if and only if $\gamma|_{M}=id_{M}$, i.e. $\gamma\in \operatorname{Gal}(L /M)$, thus $\ker \Phi=\operatorname{Gal}(L /M)$. Hence $\operatorname{Gal}(L /M)\lhd \operatorname{Gal}(L/K)$. 

    Applying the first isomorphism theorem of groups to $\Phi$, we get 
    \[
    \operatorname{Gal}(L /K)/\operatorname{Gal}(L /M) \cong \im \Phi \leq \operatorname{Gal}(M/K). 
    \]
    All those Galois groups are finite so we may count the cardinality of groups above: 
    \[
    \left| \im \Phi \right| =\frac{\left| \operatorname{Gal}(L /K) \right| }{\left|\operatorname{Gal}(L /M)\right|}=\frac{[L:K]}{[L:M]}=[M:K]=\left| \operatorname{Gal}(M /K) \right|,
    \]
    so $\im \Phi=\operatorname{Gal}(M /K)$ and 
    \[
    \operatorname{Gal}(M /K)\cong \operatorname{Gal}(L /K)/\operatorname{Gal}(L /M). 
    \]

    $(\Leftarrow)$: Suppose $\operatorname{Gal}(L /M)\lhd \operatorname{Gal}(L /K)$. Take any $\gamma\in \operatorname{Gal}(L /K)$, then $\gamma(M)$ is a subfield of $L$. Then 
    \[
    \begin{aligned}
        \operatorname{Gal}(L/ \gamma(M))&=\{\mu\in \operatorname{Gal}(L /K)|\mu(x)=x,\forall x\in \gamma(M)\}\\
        &=\{\mu\in \operatorname{Gal}(L /K)|\gamma^{-1}\mu\gamma(x)=x,\forall x\in M\}.
    \end{aligned}
    \]
    So $\mu\in \operatorname{Gal}(L /\gamma(M))$ if and only if $\gamma^{-1}\mu\gamma\in \operatorname{Gal}(L /M)$ and 
    \[
    \operatorname{Gal}(L /\gamma(M))=\gamma\operatorname{Gal}(L /M)\gamma^{-1}. 
    \]
    Since $\operatorname{Gal}(L /M)$ is normal, the equation above implies $\operatorname{Gal}(L /\gamma(M))=\operatorname{Gal}(L /M)$. Therefore 
    \[
    \gamma(M)=L^{\operatorname{Gal}(L /\gamma(M))}=L^{\operatorname{Gal}(L /M)}=M. 
    \]

Since \(
\gamma(M)=M
\) for every \(\gamma\in \operatorname{Gal}(L/K),\) we may define the restriction homomorphism
\[
\Phi:\operatorname{Gal}(L/K)\longrightarrow \operatorname{Aut}_K(M),
\qquad
\gamma\longmapsto \gamma|_M.
\]
with \(\ker\Phi
=
\{\gamma\in\operatorname{Gal}(L/K)\mid \gamma|_M=\operatorname{id}_M\}
=
\operatorname{Gal}(L/M).
\) Hence, by the First Isomorphism Theorem,
\[
|\operatorname{im}\Phi|
=
\frac{|\operatorname{Gal}(L/K)|}
     {|\operatorname{Gal}(L/M)|}.
\]
Since $L:K$ and $L:M$ are Galois,
\[
|\operatorname{Gal}(L/K)|=[L:K],
\qquad
|\operatorname{Gal}(L/M)|=[L:M].
\]
Therefore, by the Tower Law,
\[
|\operatorname{im}\Phi|
=
\frac{[L:K]}{[L:M]}
=
[M:K].
\]
By the equality established in the proof of Theorem \ref{thm-FTGT-1},
\[
[M:M^{\operatorname{im}\Phi}]=|{\operatorname{im}\Phi}|=[M:K].
\]
where \(
K\subseteq M^{\operatorname{im}\Phi}\subseteq M.
\)
Thus the Tower Law gives \(
[M^{\operatorname{im}\Phi}:K]=1
\)
and hence
\[
M^{\operatorname{im}\Phi}=K.
\]
Finally, since
\(
{\operatorname{im}\Phi}\leq \operatorname{Aut}_K(M),
\)
we have
\[
M^{\operatorname{Aut}_K(M)}
\subseteq M^{\operatorname{im}\Phi}=K.
\]
On the other hand, every $K$-automorphism fixes $K$, so
\[
K\subseteq M^{\operatorname{Aut}_K(M)}.
\]
Consequently,
\(
M^{\operatorname{Aut}_K(M)}=K
\)
and the result follows.
\end{proof}

\begin{remark}\label{rmk-GaloisExtensionEquiFixedField}
    In the proof of Theorem \ref{thm-FTGT-2}, we also proved the following result: 
    \begin{center}
        \fbox{$M :K$ is Galois extension $\iff$ $\forall \gamma\in \operatorname{Gal}(L /K), \gamma(M)=M$}
    \end{center}
\end{remark}
